\subsection{Milestone 1: Inception}
\label{sec:m1}

\subsubsection{Selection of problem and requirements engineering for ML}
\label{sec:m1-requirements}
\guidance{
  \textbf{Template:} How, and why, has the problem been selected, and what are the system requirements?\\
  \textbf{Rubric (6 pts):} How clearly does the project definition and requirements show what makes the system unique and complete?\\
  Cover functional and non-functional requirements (quality objectives, performance criteria, constraints).
  Requirements and quality objectives are typical M1 EDN entries.
}

\tbd{problem selection and system requirements}

\subsubsection{Dataset card}
\label{sec:m1-dataset-card}
\guidance{
  \textbf{Template:} Brief description of the main points of the dataset card, and link to the file in the repository.\\
  \textbf{Rubric (shared with the model card):} How well does the dataset and model documentation communicate our choices and their implications for others?\\
  Follow the Hugging Face dataset card template \cite{hf-dataset-card}.
}

\tbd{main points of the dataset card and link to it in the repository}

\subsubsection{Model card}
\label{sec:m1-model-card}
\guidance{
  \textbf{Template:} Brief description of the main points of the model card, and link to the file in the repository.\\
  Follow the Hugging Face annotated model card template \cite{hf-model-card, mitchell2019modelcards}.
  Update it in the final report once the model has changed.
}

\tbd{main points of the model card and link to it in the repository}

\subsubsection{Project coordination and communication}
\label{sec:m1-coordination}
\guidance{
  \textbf{Template:} What tools have been selected, why, and how are they used?\\
  \textbf{Rubric (4 pts):} How effectively do our coordination and communication practices support collaboration and leave a useful trace of our progress?\\
  Already in place (see \texttt{docs/docs/scrum/} and \texttt{CONTRIBUTING.md}): Scrum with GitHub Issues and a GitHub Projects board, Discord for async communication, a weekly sync in the Tuesday lab instead of a daily standup, a Definition of Done, working agreements and a sprint log, issue and PR templates.
  Show the trace: board screenshot, sprint log, linked issues and PRs.
}

\tbd{coordination and communication tools and how we use them}

\subsubsection{Selection of the cloud provider}
\label{sec:m1-cloud}
\guidance{
  \textbf{Template:} Justification of the selection (to be used in Milestone~4).\\
  The demo repo compares free tiers of AWS, Azure, GCP, Okteto and the FIB Virtech VM (free for the semester, 4\,GB RAM, 20\,GB storage); see references/course-demos.md.
  Tie the choice to our requirements (model size, memory, cost, how long it has to run).
}

\tbd{chosen cloud provider and justification}
